We provide the first annual estimates of Latvian GDP per capita for the Soviet period, 1937--95, built from physical output, transport and employment series drawn from Soviet yearbooks and archives and anchored to a 1990 purchasing-power-parity benchmark, and splice them to the existing interwar and modern series to obtain a continuous series for 1913 and 1919--2024 in international dollars. Independently constructed series agree within 3.5 per cent in every year of their overlap. The new series completes the Baltic record and permits the first national-accounts-based comparison of Latvia, Lithuania and Estonia under interwar capitalism, Soviet socialism and restored independence. Despite high social capability, none of the three converged on Western Europe before 1990: the growth of the 1930s and of 1950--68 was cancelled by the stagnation of 1968--89, and Latvia, the richest Baltic economy for most of the period from 1913 to 1988, is now the poorest. Catch-up came only after 1995, and has been conditional and uneven.
